% Requires \usepackage{booktabs}
\begin{table}[htbp]
  \centering
  \caption{Thickness-extrapolated surface stress per facet, the one L2 quantity in this chain and the input every later number depends on. tau\_xx and tau\_yy are the two in-plane axes named in the second column; the axis assignment is not conventional and the anisotropy sign is meaningless without it. POSITIVE tau IS COMPRESSIVE (CLAUDE.md sec 2): it equals MINUS the continuum surface stress f, so a facet with positive tau relaxes by EXPANDING, which was verified directly against free 2D vc-relax on all three surfaces and six axes. sigma\_res is the residual interior stress the ladder fit leaves behind, within the 0.3 kbar tolerance on all three facets. The rms column is the larger of the two per-axis fit residuals expressed as a tau, i.e. the fit's own scatter, and is not an estimate of the systematic error in the DFT. (111) is isotropic identically, not numerically: its 3-fold symmetry forbids an in-plane anisotropy. \textbullet\ Thickness definition: t = N\_C * a0\textasciicircum{}3/8 / A (bulk-equivalent, atom-counted). \textbullet\ Fit uses 8L 10L 12L 16L; 6L excluded. \textbullet\ a0 = 3.572997 A at 90/720 Ry.}
  \label{tab:t2-tau-infinity}
  \begin{tabular}{llrrrrrrl}
    \toprule
    facet & axes $x$ / $y$ & $\tau_{xx}$ & $\tau_{yy}$ & $\bar\tau$ & $\tau_{xx}-\tau_{yy}$ & $\sigma_{\rm res}$ & fit rms & level \\
    \midrule
    (100) & [110] / [1-10] & +1.032 & -5.024 & -1.996 & +6.056 & -0.283 & 0.0290 & L2 \\
    (110) & [1-10] / [001] & +2.144 & +4.397 & +3.271 & -2.253 & -0.257 & 0.0014 & L2 \\
    (111) & [1-10] / [11-2] & +0.483 & +0.483 & +0.483 & +0.000 & -0.362 & 0.0021 & L2 \\
    \bottomrule
  \end{tabular}
\end{table}
